\documentclass[a4paper,12pt]{article}
\usepackage[utf8]{inputenc}
\usepackage[T1]{fontenc}
\usepackage[french]{babel}
\usepackage{geometry}
\usepackage{xcolor}
\usepackage[explicit]{titlesec}
\usepackage{tcolorbox}
\usepackage{lmodern}
\usepackage{enumitem}
\usepackage{booktabs}
\usepackage{tabularx}
\usepackage{array}
\usepackage{fancyhdr}
\usepackage{hyperref}
\usepackage{graphicx}
\usepackage{float}

\geometry{margin=2.5cm, top=3cm, bottom=3cm}

% ── Couleurs ──────────────────────────────────────────────────────────────────
\definecolor{primary}{RGB}{0, 70, 140}
\definecolor{secondary}{RGB}{230, 240, 255}
\definecolor{lightgray}{RGB}{245, 245, 245}
\definecolor{darkgray}{RGB}{80, 80, 80}

% ── En-tête / Pied de page ────────────────────────────────────────────────────
\pagestyle{fancy}
\fancyhf{}
\lhead{\textcolor{darkgray}{\small BTS SIO — SISR | AP-2}}
\rhead{\textcolor{darkgray}{\small Projet ASSIZ — M2L}}
\cfoot{\textcolor{darkgray}{\thepage}}
\renewcommand{\headrulewidth}{0.4pt}

% ── Style des sections ────────────────────────────────────────────────────────
\titleformat{\section}
  {\normalfont\Large\bfseries\sffamily\color{primary}}
  {}
  {0pt}
  {\colorbox{secondary}{\parbox{\dimexpr\textwidth-2\fboxsep}{\thesection\quad #1}}}

\titleformat{\subsection}
  {\normalfont\normalsize\bfseries\sffamily\color{primary}}
  {\thesubsection}
  {0.5em}
  {#1}
  [\vspace{-0.3em}\textcolor{primary!40}{\hrule height 0.5pt}\vspace{0.3em}]

% ── Style des listes ──────────────────────────────────────────────────────────
\setlist[itemize,1]{label=\textcolor{primary}{$\bullet$}, leftmargin=2em, labelsep=1em}
\setlist[itemize,2]{label=\textcolor{primary!60}{--}, leftmargin=2em, labelsep=0.8em}
\setlist[enumerate,1]{label=\textcolor{primary}{\arabic*.}, leftmargin=2em, labelsep=1em}
\setlist{itemsep=2pt, parsep=2pt}

% ── Boîte info ────────────────────────────────────────────────────────────────
\tcbuselibrary{skins}
\newtcolorbox{infobox}{
  colback=secondary, colframe=primary,
  boxrule=0.5pt, leftrule=4pt,
  sharp corners, left=8pt, right=8pt, top=6pt, bottom=6pt
}

% ── Tableau stylé ─────────────────────────────────────────────────────────────
\newcolumntype{L}[1]{>{\raggedright\arraybackslash}p{#1}}
\newcolumntype{C}[1]{>{\centering\arraybackslash}p{#1}}

\hypersetup{colorlinks=true, linkcolor=primary, urlcolor=primary, citecolor=primary}

% =============================================================================
\begin{document}
\thispagestyle{empty}

% ── Page de garde ─────────────────────────────────────────────────────────────
\begin{center}
  \vspace*{2cm}
  {\color{primary}\rule{\textwidth}{2pt}}
  \vspace{0.8em}

  {\Huge\bfseries\sffamily\color{primary} Projet ASSIZ}\\[0.4em]
  {\Large\sffamily\color{darkgray} Couverture Wifi des Assises de l'Escrime}\\[0.2em]
  {\normalsize\color{darkgray} Maison des Ligues de Lorraine — AP-2 SISR}

  \vspace{0.8em}
  {\color{primary}\rule{\textwidth}{2pt}}

  \vspace{2cm}

  \begin{tcolorbox}[colback=lightgray, colframe=primary!30, sharp corners, width=0.7\textwidth]
    \centering
    \textbf{Groupe :} Maël — Joan — Camille — Anna\\[0.3em]
    \textbf{Formation :} BTS SIO option SISR\\[0.3em]
    \textbf{Année :} 2025–2026\\[0.3em]
    \textbf{Date :} \today
  \end{tcolorbox}
\end{center}

\vfill
{\small\color{darkgray}\centering Réseau CERTA — Contexte M2L\par}

\newpage
\tableofcontents
\newpage

% =============================================================================
\section{Contexte et problématique}

\subsection{Présentation du contexte}

% Décrire la situation de départ : la ligue d'Escrime organise une manifestation
% sous chapiteaux (2500 visiteurs), à quelques centaines de mètres de la M2L.
% Le service informatique doit prototyper l'infrastructure réseau nomade.

\subsection{Problématique identifiée}

% Expliquer le besoin : fournir un accès Wifi sécurisé et séparé
% pour les organisateurs et les visiteurs, sans que le public
% puisse accéder aux ressources de l'orga.

\subsection{Objectifs du projet}

\begin{itemize}
  \item
  \item
  \item
\end{itemize}

% =============================================================================
\section{Infrastructure réseau}

\subsection{Matériel réseau}

\begin{itemize}
  \item \textbf{Routeur ISR4321} — Arrivée internet
  \item \textbf{Routeur ISR4331} — Gestion du réseau local
  \item \textbf{Switch 2960-24TT} — Switch principal, connecte les machines du réseau
  \item \textbf{AP-PT} — Access Point, gère le réseau Wi-Fi (×5 pour couvrir toute la zone)
\end{itemize}

\subsection{Matériel informatique}

\begin{itemize}
  \item \textbf{PC-PT — PC-1-Accueil} — Poste de travail pour l'accueil, connecté au switch
  \item \textbf{PC-PT — PC-2-Orga} — Poste de travail pour l'organisation, connecté au switch
  \item \textbf{Laptop-PT — Portable-Orga} — Ordinateur portable pour l'organisation, connecté à l'AP
  \item \textbf{Laptop-PT — Portable-Visiteur} — Ordinateur portable dédié aux visiteurs, connecté à l'AP
\end{itemize}

\subsection{Organisation du réseau}

\begin{itemize}
  \item Le routeur ISR4321 est connecté à Internet et gère la connexion entrante.
  \item Le routeur ISR4331 est connecté au routeur ISR4321 et gère le réseau local, y compris les VLANs et les règles de sécurité.
  \item Le switch 2960-24TT est connecté au routeur ISR4331 et connecte les postes de travail et les APs.
  \item Les AP-PT sont connectés au switch et fournissent une couverture Wi-Fi pour les zones d'accueil et d'organisation.
  \item Les postes fixes sont connectés au switch ; les portables se connectent via les AP.
\end{itemize}

\subsection{Schéma réseau}

% Insérer ici le schéma de topologie réseau (export Packet Tracer ou Visio)
% \begin{figure}[H]
%   \centering
%   \includegraphics[width=0.9\textwidth]{schema_assiz.png}
%   \caption{Topologie réseau — Projet ASSIZ}
% \end{figure}

\begin{infobox}
  \textit{Schéma à insérer ici — export depuis Cisco Packet Tracer ou Draw.io}
\end{infobox}

\subsection{Plan d'adressage IP}

\begin{tabularx}{\textwidth}{L{3.5cm} C{2.5cm} C{2.5cm} C{2.5cm} L{3cm}}
  \toprule
  \textbf{Équipement} & \textbf{VLAN} & \textbf{Adresse IP} & \textbf{Masque} & \textbf{Passerelle} \\
  \midrule
  Routeur ISR4331 (LAN) & — & & & — \\
  Switch 2960-24TT & 1 & & & \\
  PC-1-Accueil & 2 & & & \\
  PC-2-Orga & 2 & & & \\
  Portable-Orga & 2 & DHCP & & \\
  Portable-Visiteur & 3 & DHCP & & \\
  AP-PT (×5) & 1 & & & \\
  \bottomrule
\end{tabularx}

\subsection{VLANs}

\begin{tabularx}{\textwidth}{C{2cm} C{3cm} X C{3.5cm}}
  \toprule
  \textbf{VLAN ID} & \textbf{Nom} & \textbf{Équipements} & \textbf{Plage réseau} \\
  \midrule
  VLAN 1 & Management & Switch, APs & \\
  VLAN 2 & Orga & PC-1-Accueil, PC-2-Orga, Portable-Orga & 10.10.10.0/24 \\
  VLAN 3 & Public & Portable-Visiteur & 172.31.1.0/24 \\
  \bottomrule
\end{tabularx}

% =============================================================================
\section{Configuration du réseau}

\subsection{Configuration des routeurs}

\subsubsection{Routeur ISR4321 — Connexion Internet}

% Décrire et coller la configuration : interface WAN, NAT outside, route par défaut

\begin{tcolorbox}[colback=black!90, coltext=green!70!white, fontupper=\ttfamily\small,
  title={\color{white}\small Configuration ISR4321}, colbacktitle=primary]
% Coller ici les commandes IOS
\end{tcolorbox}

\subsubsection{Routeur ISR4331 — Réseau local}

% Décrire et coller la configuration : sous-interfaces (router-on-a-stick),
% encapsulation 802.1Q, adresses IP par VLAN

\begin{tcolorbox}[colback=black!90, coltext=green!70!white, fontupper=\ttfamily\small,
  title={\color{white}\small Configuration ISR4331}, colbacktitle=primary]
% Coller ici les commandes IOS
\end{tcolorbox}

\subsection{Configuration NAT}

% Expliquer et coller la configuration NAT (ip nat inside / outside,
% access-list, ip nat inside source)

\begin{tcolorbox}[colback=black!90, coltext=green!70!white, fontupper=\ttfamily\small,
  title={\color{white}\small Configuration NAT}, colbacktitle=primary]
% Coller ici les commandes IOS
\end{tcolorbox}

\subsection{Configuration DHCP}

% 2 pools DHCP : un pour le VLAN orga (10.10.10.0/24) et un pour le VLAN public (172.31.1.0/24)

\begin{tcolorbox}[colback=black!90, coltext=green!70!white, fontupper=\ttfamily\small,
  title={\color{white}\small Configuration DHCP}, colbacktitle=primary]
% Coller ici les commandes IOS
\end{tcolorbox}

\subsection{Configuration du switch}

% Création des VLANs, ports en mode access (vers postes fixes),
% ports en mode trunk (vers routeur et APs)

\begin{tcolorbox}[colback=black!90, coltext=green!70!white, fontupper=\ttfamily\small,
  title={\color{white}\small Configuration Switch 2960}, colbacktitle=primary]
% Coller ici les commandes IOS
\end{tcolorbox}

\subsection{Configuration des Access Points}

% Pour chaque AP : 2 SSIDs (public diffusé WPA2, orga non diffusé WPA2),
% association SSID → VLAN

\subsubsection{SSID "orga"}

\begin{itemize}
  \item Diffusion SSID : non diffusé (hidden)
  \item Sécurité : WPA2
  \item Clé : (connue des organisateurs uniquement)
  \item VLAN associé : VLAN 2
\end{itemize}

\subsubsection{SSID "public"}

\begin{itemize}
  \item Diffusion SSID : diffusé
  \item Sécurité : WPA2
  \item Clé : (indiquée sur le badge visiteur)
  \item VLAN associé : VLAN 3
  \item Isolation entre postes : activée
\end{itemize}

\subsection{Règles de filtrage}

% Décrire les ACL ou règles de pare-feu empêchant le VLAN public
% d'accéder aux ressources du VLAN orga (imprimante, traceur, afficheur géant)

\begin{tcolorbox}[colback=black!90, coltext=green!70!white, fontupper=\ttfamily\small,
  title={\color{white}\small Règles de filtrage (ACL)}, colbacktitle=primary]
% Coller ici les commandes IOS
\end{tcolorbox}

% =============================================================================
\section{Tests et validation}

\subsection{Plan de tests}

\begin{tabularx}{\textwidth}{C{0.5cm} L{4.5cm} L{4cm} C{2cm} C{2cm}}
  \toprule
  \textbf{\#} & \textbf{Test} & \textbf{Action} & \textbf{Résultat attendu} & \textbf{Résultat obtenu} \\
  \midrule
  1 & Attribution DHCP VLAN orga & Connecter Portable-Orga sur SSID "orga" & IP en 10.10.10.x & \\
  2 & Attribution DHCP VLAN public & Connecter Portable-Visiteur sur SSID "public" & IP en 172.31.1.x & \\
  3 & Accès Internet depuis VLAN orga & Ping 8.8.8.8 depuis Portable-Orga & Succès & \\
  4 & Accès Internet depuis VLAN public & Ping 8.8.8.8 depuis Portable-Visiteur & Succès & \\
  5 & Isolation VLAN public → orga & Ping PC-2-Orga depuis Portable-Visiteur & Échec (filtré) & \\
  6 & Communication interne VLAN orga & Ping PC-1-Accueil depuis Portable-Orga & Succès & \\
  \bottomrule
\end{tabularx}

\subsection{Captures d'écran}

% Insérer les captures : attribution DHCP, résultats de ping, interface AP, etc.

\begin{infobox}
  \textit{Captures d'écran à insérer ici (DHCP, pings, interface de configuration AP)}
\end{infobox}

% =============================================================================
\section{Sauvegarde de la configuration (TFTP)}

\subsection{Procédure de sauvegarde}

% Décrire la procédure : adresse du serveur TFTP, commande copy running-config tftp

\begin{tcolorbox}[colback=black!90, coltext=green!70!white, fontupper=\ttfamily\small,
  title={\color{white}\small Sauvegarde vers TFTP}, colbacktitle=primary]
% copy running-config tftp
% Address or name of remote host []?
% Destination filename []?
\end{tcolorbox}

\subsection{Vérification}

% Confirmer que le fichier est bien présent sur le serveur TFTP

% =============================================================================
\section{Avantages et limites de la solution}

\begin{tabularx}{\textwidth}{X X}
  \toprule
  \textbf{Avantages} & \textbf{Limites / Inconvénients} \\
  \midrule
  & \\
  & \\
  & \\
  & \\
  \bottomrule
\end{tabularx}

% =============================================================================
\section*{Conclusion}
\addcontentsline{toc}{section}{Conclusion}

% Résumer ce qui a été mis en place, les compétences mobilisées,
% et ce que ce projet apporte dans le cadre du référentiel SISR2.

\end{document}
